\section{BV-BRST construction}
\label{sec:BV_construction}
\label{ch:bv-construction}

The bridge theorem (Theorem~\ref{thm:physics-bridge}) identifies the
BV-BRST data of a 3d HT theory on $\C_z \times \R_t$ with a
logarithmic $\SCchtop$-algebra under the product HT gauge,
factorized logarithmic parametrix, one-loop finiteness, and
product-polynomial interaction hypotheses. The prefactorization algebra
$U \mapsto \mathsf{Obs}(U)$ produces the open factorization category;
the quantum master equation
\begin{equation}
\label{eq:QME-bv-construction}
\hbar\,\Delta_{\mathrm{BV}}\, S_{\mathrm{eff}} + \tfrac{1}{2}\{S_{\mathrm{eff}}, S_{\mathrm{eff}}\}_{\mathrm{BV}} = 0
\end{equation}
is the open/closed Maurer--Cartan equation.

\begin{remark}[The $\hbar$ convention and the quantum-group bridge]
\label{rem:bv-hbar-convention}
$\hbar$ is the formal expansion parameter of the BV
quantization with $|\hbar|=0$; $\hbar\Delta_{\mathrm{BV}}$ adds the loop
expansion to the classical antibracket $\tfrac12\{S,S\}$. This is the Costello
convention~\cite{CG17}: the QME~\eqref{eq:QME-bv-construction} reduces at
$\hbar=0$ to the classical master equation $\{S,S\}=0$. The spectral parameter $\hbar$ in
$R(z) = \exp(r(z)) = 1 + \hbar\,r(z) + O(\hbar^2)$ is the BV $\hbar$ when both are computed from the same boundary
chiral algebra. The BV $\hbar$ carries no $\pi i$ factor; the $2\pi i$
convention enters only in residue evaluations of weight forms over
$\FM_k(\C)$ via $\frac{1}{2\pi i}\oint$.
\end{remark}

\begin{remark}[Bar complex and BV complex: exact scope]
\label{rem:bar-equals-BV-bridge}
For the boundary chiral algebra $A_\partial$ of a 3d HT theory satisfying
Theorem~\ref{thm:physics-bridge}, the chiral bar complex
$\barB^{\mathrm{ch}}(A_\partial) = T^c(s^{-1}\bar A_\partial)$
compares with the BV-BRST complex
$\Obs_{\mathrm{bulk}} = (\widehat{\Sym}(\mathcal E^\vee), Q_{\mathrm{BRST}})$
through a comparison datum constructed below. At
genus zero the chain comparison is the Costello--Gwilliam /
Drinfeld--Sokolov--Hochschild bridge. In higher genus the comparison is coderived; a strict equality between
$d_{\overline B}$ and
$Q_{\mathrm{BRST}}=\{S_{\mathrm{eff}},-\}_{\mathrm{BV}}
+\hbar\Delta_{\mathrm{BV}}$ requires the
diagonal-residue BV operator to exist, be second order, square to zero,
and be functorial. The comparison is the
chain-level form of the universal-holography master theorem
(Theorem~\ref{thm:universal-holography-master}).
The Koszul conductor $K = -c_{\mathrm{ghost}}(\mathrm{BRST})$ from the
Volume~I universal-conductor theorem equals the level shift controlling
the QME at one loop: $\kappaChHodge(A_\partial) + \kappaChHodge(A_\partial^!) = K$ in the
scalar channel, with the ghost central charge absorbing the matter
complementarity. The QME \eqref{eq:QME-bv-construction} is the
chain-level incarnation of $D_A\Theta_A + \tfrac12[\Theta_A, \Theta_A] = 0$
on the bar side: open (BV) and closed (modular Maurer--Cartan) faces of
the same Maurer--Cartan datum.
\end{remark}

\begin{remark}[BV $=$ bar scope: class G/L proved, class C/M conditional]
\label{rem:bv-equals-bar-scope}
The chain-level identification $(\Obs_{\mathrm{bulk}}, Q_{\mathrm{BRST}})
\simeq (\barB^{\mathrm{ch}}(A_\partial), d_{\overline B})$ holds
unconditionally for the abelian--Sugawara class~G (Heisenberg, lattice)
and the affine Kac--Moody class~L at non-critical level via the explicit
homotopy transfer~\cite{CG17}*{Theorem~5.2.2.1} of
\S\ref{subsec:analytic} Step~4. For the symplectic bosons /
$\beta\gamma$ class~C and the Virasoro / $W_N$ class~M, the
identification depends on the harmonic component of the BV propagator.
For class~C the required input is the vanishing, with compatible lifts,
of the obstruction classes
\(\mathfrak h_{g,r}^{\mathsf C}\) of
Definition~\ref{def:class-c-harmonic-decoupling-obstruction}; this is
established only at the printed genus-one quartic test case. For
class~M the raw direct-sum chain-level statement is false, and the
replacement is the weight-completed comparison below. The cohomological identification
$H^\bullet(\Obs_{\mathrm{bulk}}) \simeq
H^\bullet(\barB^{\mathrm{ch}}(A_\partial))$ holds for all four classes
via the Drinfeld--Sokolov/Hochschild bridge
(Theorem~\ref{thm:chd-ds-hochschild}); the chain-level upgrade for class
C is conditional on the classes \(\mathfrak h_{g,r}^{\mathsf C}\),
while class~M closes only in the weight-completed category by
Proposition~\ref{prop:bv-bar-class-m-weight-completed}. The
QME~\eqref{eq:QME-bv-construction} is the open form;
$D_A\Theta_A + \tfrac12[\Theta_A, \Theta_A] = 0$ is the closed form. Both
present the same Maurer--Cartan element under the bar/BV identification on
class G/L, on class~C after vanishing of
\(\mathfrak h_{g,r}^{\mathsf C}\), and on class~M in the
weight-completed category.
\end{remark}

\subsection{Field content, action, and BV data}
\label{subsec:BV-data}

Fix a cyclic dg Lie algebra \((\mathfrak a,d_{\mathfrak a},[-,-],
\langle-,-\rangle_{\mathfrak a})\). The local 3d
holomorphic--topological bridge starts from the BV datum
\[
(\mathcal E,Q,\omega_{\mathrm{BV}},I,B)
\]
on \(\R_t\times\C_z\), where
\begin{equation}\label{eq:3dht-fields}
 \mathcal E
 =
 \Omega^\bullet(\R_t)\widehat\otimes
 \Omega^{0,\bullet}(\C_z)\otimes\mathfrak a[1],
 \qquad
 Q=d_t+\dbar_z+d_{\mathfrak a}.
\end{equation}
The BV pairing is
\begin{equation}\label{eq:3dht-bv-pairing}
 \omega_{\mathrm{BV}}(\alpha,\beta)
 =
 \int_{\R\times\C} dz\wedge
 \langle\alpha,\beta\rangle_{\mathfrak a},
\end{equation}
where the \(dt\) and \(d\bar z\) components are supplied by the form
degrees of \(\alpha,\beta\). The interaction and classical action are
\begin{align}
 I(A)
 &=
 \frac16\int_{\R\times\C}dz\wedge
 \langle A,[A,A]\rangle_{\mathfrak a},\label{eq:3dht-interaction}\\
 S_{\mathrm{cl}}[A]
 &=
 \int_{\R\times\C}dz\wedge
 \left(\frac12\langle A,QA\rangle_{\mathfrak a}
 +\frac16\langle A,[A,A]\rangle_{\mathfrak a}\right).
 \label{eq:3dht-classical-action}
\end{align}
The boundary condition \(B\) is a \(Q\)-compatible Lagrangian
subcomplex along \(t=0\), so the BV boundary term vanishes. Invariance
of \(\langle-,-\rangle_{\mathfrak a}\), \(Q^2=0\), and the Jacobi
identity give the classical master equation
\begin{equation}\label{eq:3dht-cme}
 \{S_{\mathrm{cl}},S_{\mathrm{cl}}\}_{\mathrm{BV}}=0.
\end{equation}

Choose a gauge-fixing adjoint \(Q^*\). Let \(H=[Q,Q^*]\) be the heat
operator, \(K_s=e^{-sH}\) its kernel, and define the scale propagator
\begin{equation}\label{eq:3dht-heat-propagator}
 P_{\epsilon<L}
 =
 \int_{\epsilon}^{L} Q^*K_s\,ds,
 \qquad
 Q_{\boxtimes}P_{\epsilon<L}=K_L-K_\epsilon.
\end{equation}
Here \(P_{\epsilon<L}=P_{\epsilon<L}(x,y)\) and \(K_s=K_s(x,y)\) are
distributional kernels on \((\R\times\C)^2\), and
\[
 Q_{\boxtimes}=Q_x\boxtimes 1+1\boxtimes Q_y
\]
is the differential induced by \(Q\) on kernels, with the Koszul sign on
the second factor determined by the BV pairing. Equivalently, as an
operator on compactly supported fields,
\begin{equation}\label{eq:3dht-propagator-kernel-homotopy}
 [Q,P_{\epsilon<L}]=K_L-K_\epsilon .
\end{equation}
After local counterterms have removed the \(\epsilon\to0\) singular
part, write \(P_{0<L}\) for the renormalized limit. The effective
interaction at scale \(L\) is the Costello graph sum
\begin{equation}\label{eq:3dht-effective-interaction}
 I[L]
 =
 \sum_{\Gamma}
 \frac{\hbar^{b_1(\Gamma)}}{|\operatorname{Aut}\Gamma|}
 W_\Gamma(P_{0<L},I),
\end{equation}
with graph weight
\begin{equation}\label{eq:3dht-graph-weight}
 W_\Gamma(P,I)
 =
 \int_{\Conf_{V(\Gamma)}(\R\times\C)}
 \bigwedge_{e\in E(\Gamma)}P_e\cdot
 \bigotimes_{v\in V(\Gamma)} I_v.
\end{equation}
The scale-\(L\) quantum master equation is
\begin{equation}\label{eq:3dht-scale-qme}
 QI[L]+\hbar\Delta_L I[L]+\frac12\{I[L],I[L]\}_L=0.
\end{equation}
Boundary \(A_\infty\)-chiral operations are extracted by placing the
external legs on \(B\):
\begin{equation}\label{eq:3dht-boundary-operations}
 m_k^{\mathrm{BV}}
 =
 \sum_{\substack{\Gamma\ \mathrm{connected}\\
 k\ \mathrm{external\ boundary\ legs}\\
 1\ \mathrm{output}}}
 \frac{\hbar^{b_1(\Gamma)}}{|\operatorname{Aut}\Gamma|}
 \operatorname{Res}_{\partial\FM_\Gamma} W_\Gamma(P_{0<L},I).
\end{equation}
The comparison with the bar construction is not an additional identity
between two unrelated formulas.  It is the transfer datum identifying
the boundary observables with the chiral bar model.  More precisely,
let \(A_\partial\) be the boundary chiral algebra obtained from \(B\)
by \(\eta_A^\partial\), and let \(D_A\) be the degree-\(1\)
coderivation of
\(\barB^{\mathrm{ch}}(A_\partial)=T^c(s^{-1}\bar A_\partial)\), with
\[
 m_k^{\mathrm{bar}}
 =
 \pi_1D_A\big|_{(s^{-1}A_\partial)^{\otimes k}} .
\]
The bridge includes a boundary strong deformation retract, compatible
with the residue kernels, with maps
\[
 i_\partial\colon (A_\partial,m_1)\to
 \bigl(\Obs^\partial(\mathcal E,B),Q_{\mathrm{BRST}}^L\bigr),
 \qquad
 p_\partial\colon
 \bigl(\Obs^\partial(\mathcal E,B),Q_{\mathrm{BRST}}^L\bigr)
 \to (A_\partial,m_1),
\]
and homotopy \(h_\partial\) satisfying
\[
 h_\partial Q_{\mathrm{BRST}}^L+Q_{\mathrm{BRST}}^Lh_\partial
 =\mathrm{id}-i_\partial p_\partial .
\]
Homological perturbation transfers the graph operations
\eqref{eq:3dht-boundary-operations} to a coderivation
\(D_A^{\mathrm{BV}}\) on \(T^c(s^{-1}\bar A_\partial)\).  The rigorous
comparison statement is the transferred equality
\begin{equation}\label{eq:3dht-bv-bar-transfer-comparison}
 \pi_1D_A^{\mathrm{BV}}\big|_{(s^{-1}A_\partial)^{\otimes k}}
 =
 \pi_1D_A\big|_{(s^{-1}A_\partial)^{\otimes k}} .
\end{equation}
Equivalently, the raw residue operations \(m_k^{\mathrm{BV}}\) are
identified with \(m_k^{\mathrm{bar}}\) by the \(A_\infty\)
quasi-isomorphism induced by \(p_\partial,i_\partial,h_\partial\);
strict equality is asserted only in the transferred boundary model.

\begin{proposition}[Standard-family BV realisation data;
\ClaimStatusConditional; licensing $\alpha+\beta+\gamma+\delta$]
\label{prop:standard-family-bv-realisation-data}
For every standard family used in the volume, the physical input
\(\Xi_A\) is not a name for a hoped-for theory.  It consists of a
field complex \(\mathcal E_A\), BRST operator \(Q_A\), interaction
\(I_A\), BV pairing \(\omega_{\mathrm{BV},A}\), gauge-fixing
operator \(Q_A^*\), boundary Lagrangian \(B_A\), scale-\(L\) solution
of the QME, and comparison maps
\[
 \eta_A^\partial\colon \Obs^\partial(\mathcal E_A,B_A)\longrightarrow A,
 \qquad
 \chi_A\colon \Zderch{A}\longrightarrow
 H^\bullet\Obs^{\mathrm{bulk}}(\mathcal E_A).
\]
The following table records the rows actually licensed here.
\begin{center}
\small
\begin{tabular}{p{0.16\textwidth}p{0.24\textwidth}p{0.25\textwidth}p{0.24\textwidth}}
family & BV source & boundary Lagrangian and QME &
comparison maps \\
\hline
Heisenberg/lattice &
abelian cyclic complex of currents; \(I_A=0\) or quadratic CS term &
Neumann/Fock Lagrangian; Gaussian QME and heat-kernel gauge fixing &
\(\eta^\partial_A\) is the current OPE; \(\chi_A\) is the free chiral Hochschild comparison \\
affine \(V_k(\fg)\) &
mixed BF/Chern--Simons fields \((\alpha,\beta)\) of Proposition~\ref{prop:affine-mixed-bf-dictionary} &
\(i^*\alpha=0\) chiral boundary Lagrangian; one-loop-exact QME by Theorem~\ref{thm:affine-half-space-bv} &
\(\eta^\partial_A\) is Corollary~\ref{cor:boundary-affine-voa}; \(\chi_A\) is the affine bar/BV comparison \\
\(\beta\gamma\)/symplectic boson &
shifted-cotangent free-field complex with finite contact interaction &
standard free boundary Lagrangian; QME after vanishing/lifting of \(\mathfrak h_{g,r}^{\mathsf C}\) and contact estimates &
\(\eta^\partial_A,\chi_A\) require the vanishing/lifting of \(\mathfrak h_{g,r}^{\mathsf C}\) \\
\(\mathrm{Vir}_c,\mathcal W_N\) &
principal DS BRST reduction of the affine row, not an independent bare BV source &
DS ghost-zero boundary gauge; QME transported under \(\hypKZSDR+\hypStokes+\hypTLift\) &
\(\eta^\partial_A,\chi_A\) are the DS--Hochschild/HPL comparison maps in the weight-completed ambient
\end{tabular}
\end{center}
Thus a statement about a standard-family HT realisation is licensed
only after the row supplies these data; the algebra \(A\) alone never
constructs its BV theory.
\end{proposition}

\begin{definition}[Local observables]
\label{def:local-obs}
For an open $U\subset \R\times\C$, define $\mathsf{Obs}(U)$ as the completed symmetric algebra $\widehat{\operatorname{Sym}}(\mathcal{E}(U)^\vee)$ of distributional sections of the dual field bundle, equipped with the scale-\(L\) BV differential \(Q_{\mathrm{BRST}}^L\) and the BV antibracket. The renormalized time-ordered product $\star_U$ is built from the scale propagator \(P_{0<L}\) and the interaction vertices \(I_v\) via the standard perturbative expansion (see e.g.\ \cite{CG17}).
\end{definition}

\subsection{Proof of Theorem~\ref{thm:physics-bridge}}
\label{subsec:analytic}

\begin{proposition}[Prefactorization structure under the bridge hypotheses]
\label{prop:prefact-Obs}
Under the conditions of Theorem~\ref{thm:physics-bridge}, $U\mapsto \mathsf{Obs}(U)$ is a logarithmic $\SCchtop$-algebra (Definition~\ref{def:log-SC-algebra}).
\end{proposition}

\begin{proof}

\medskip\noindent\textbf{Step 1: Classical master equation.}
The datum \((\mathcal E,Q,\omega_{\mathrm{BV}},I,B)\) of
\eqref{eq:3dht-fields}--\eqref{eq:3dht-cme} gives a
\((-1)\)-symplectic graded space of fields on \(\R\times\C\) and a
local action functional \(S_{\mathrm{cl}}\). The quadratic part has
Hamiltonian vector field \(Q\), and \(Q^2=0\). The cubic term is
Hamiltonian for the dg Lie bracket on \(\mathfrak a\); invariance of
\(\langle-,-\rangle_{\mathfrak a}\) and Jacobi give
\(\{S_{\mathrm{cl}},S_{\mathrm{cl}}\}_{\mathrm{BV}}=0\). The boundary
condition \(B\) is Lagrangian and \(Q\)-compatible, so integration by
parts produces no boundary anomaly at \(t=0\).

Quantisation replaces \(S_{\mathrm{cl}}\) by the scale-\(L\)
interaction \(I[L]\) of \eqref{eq:3dht-effective-interaction}. The
regularised BV differential is
\[
 Q_{\mathrm{BRST}}^L
 =
 Q+\{I[L],-\}_L+\hbar\Delta_L .
\]
The scale quantum master equation \eqref{eq:3dht-scale-qme} is exactly
\((Q_{\mathrm{BRST}}^L)^2=0\). One-loop finiteness~\cite{GRW21}
ensures that the \(\epsilon\to0\) limit defining \(P_{0<L}\) and
\(\Delta_L\) exists after local counterterms. For each open
\(U\subset\R\times\C\), the cochain complex
\((\Obs(U),Q_{\mathrm{BRST}}^L)\) is therefore well-defined.

\medskip\noindent\textbf{Step 2: Propagator.}
The gauge-fixing adjoint \(Q^*\) gives the heat kernel \(K_s=e^{-sH}\)
and propagator \(P_{\epsilon<L}=\int_\epsilon^L Q^*K_s\,ds\) of
\eqref{eq:3dht-heat-propagator}, satisfying the kernel homotopy identity
\eqref{eq:3dht-propagator-kernel-homotopy}. Its singular product-HT
part is the Cauchy/step kernel
\[
 P_{\mathrm{sing}}(z,t)=\frac{\Theta(t)}{2\pi z}
 \otimes\langle-,-\rangle_{\mathfrak a}^{-1},
\]
up to a smooth \(Q\)-exact correction. Meromorphicity in \(z\) with a
simple pole at \(z_1=z_2\) ensures that correlators are rational
functions of the holomorphic insertion points, with singularities
governed by the OPE. For tree-level diagrams, rationality follows from
iterating the Cauchy integral; for loop diagrams, the one-loop
finiteness condition controls the coincident-point limit, and Hartogs
extension gives rationality away from collision loci. The heat-kernel
decay in the topological direction gives absolute convergence of
ordered-interval integrals in \(\R\). Translation invariance guarantees
that the resulting operations depend only on relative positions.

\medskip\noindent\textbf{Step 3: Factorization structure maps.}
Let $U_1,\ldots,U_n \subset \R\times\C$ be pairwise disjoint opens contained in a larger open $V$. The factorization product
\[
 \mu_{U_1,\ldots,U_n}^V \;:\;
 \Obs(U_1)\otimes \cdots \otimes \Obs(U_n)
 \;\longrightarrow\;
 \Obs(V)
\]
is defined by the renormalized time-ordered product: given observables ${\mathcal O}_i\in\Obs(U_i)$ supported in $U_i$, the product $\mu({\mathcal O}_1\otimes\cdots\otimes{\mathcal O}_n)$ is computed by the graph weights \eqref{eq:3dht-graph-weight}, with vertices at the insertion points and internal edges weighted by \(P_{0<L}\). Since the $U_i$ are disjoint, propagators connecting different opens have non-coincident holomorphic arguments $z_j\ne z_k$, so the meromorphic singularity in \(P_{\mathrm{sing}}\) is never encountered, and the heat-kernel decay in \(t\) gives convergence of the integrals over the topological separations. The BRST compatibility $Q_{\mathrm{BRST}}^L\circ\mu = \mu\circ (Q_{\mathrm{BRST}}^L)^{\otimes n}$ follows from the scale quantum master equation \eqref{eq:3dht-scale-qme}.

\medskip\noindent\textbf{Step 4: Convergence via FM compactification.}
When insertions approach coincidence, $\FM_k(\C)$ resolves the singularities. The weight forms extend to logarithmic forms $\omega_k^{\mathrm{hol}}\in\Omega^\bullet_{\log}(\FM_k(\C))$, which are integrable on the compact space $\FM_k(\C)$. The factored structure $\omega_k = \omega_k^{\mathrm{hol}}\wedge\omega_k^{\mathrm{top}}$ with $\omega_k^{\mathrm{top}}\in C_\bullet(\Conf_k(\R))$ means the integrals converge separately in each direction. By Theorem~\ref{thm:Arnold_full_proof} (Appendix~\ref{app:FM_Stokes}), codimension-2 corner contributions cancel via the Arnold--Orlik--Solomon relations, so Stokes' theorem on $\FM_k(\C)$ (a manifold with corners) yields well-defined renormalized products. Any ambiguity in the renormalization scheme is \(Q_{\mathrm{BRST}}^L\)-exact: two choices of propagator \(P,P'\) with the same singular part and boundary condition differ by a smooth kernel \(R=P-P'\), and the resulting change in each \(m_k^{\mathrm{BV}}\) is a \(Q_{\mathrm{BRST}}^L\)-coboundary by the homotopy transfer formula \cite[Theorem~5.2.2.1]{CG17}. The quasi-isomorphism class of \((\Obs(U), Q_{\mathrm{BRST}}^L)\) is therefore independent of the renormalization scheme.

\medskip\noindent\textbf{Step 5: Swiss-cheese structure.}
The factorization products assemble into the two-colored operad $\SCchtop$. The holomorphic insertions are governed by $\FM_k(\C)$: the boundary strata of the FM compactification encode the operadic composition maps $\gamma_{\mathrm{ch}}$ of the closed color, recording how clusters of holomorphic points collide. The topological insertions are governed by $E_1(m) \cong \Conf_m(\R)$: the ordered-interval structure of configurations on $\R$ provides the operadic composition $\gamma_{\mathrm{top}}$ of the open color. The factorization maps $\mu^V_{U_1,\ldots,U_n}$ on rectangles $D_i\times I_j$ (disks in $\C$, intervals in $\R$) are therefore controlled by the insertion maps of $\FM_k(\C)\times E_1(m)$, the spaces of operations of the depth-zero associated graded $\gr\,\SCchtop$ (Proposition~\ref{prop:gr-chiral}); the rectangle-basis construction is a gr-level (decoupled) statement, and the coupled bulk-to-boundary strata of $\SCchtop$ enter only through boundary-incidence limits not taken here. The no-open-to-closed constraint is automatic: by Definition~\ref{def:SC-operations}(i), $\SCchtop(\ldots,\mathrm{top},\ldots;\,\mathrm{ch})=\varnothing$ whenever any input carries the open color. Concretely, a closed-output operation requires all inputs to be holomorphic insertions parametrized by $\FM_k(\C)$; no topological (interval-ordered) input can appear. This encodes the physical directionality that bulk observables restrict to the boundary but boundary data does not extend to the bulk.

\medskip\noindent\textbf{Conclusion.}
The assignment $U\mapsto (\Obs(U), Q_{\mathrm{BRST}}^L)$ with the factorization products $\mu^V_{U_1,\ldots,U_n}$ satisfies the axioms of a holomorphic--topological prefactorization algebra: isotopy invariance along $\R$ follows from local constancy (the singular propagator depends on $t$ through the step kernel, with heat-kernel smoothing terms), and holomorphic dependence on $\C$-positions follows from the meromorphicity of \(P_{\mathrm{sing}}\). On rectangles $D\times I$ the structure is locally constant, as claimed.
\end{proof}

\begin{remark}[Lagrangian interpretation]
\label{rem:lagrangian-geometry}
The product decomposition of Definition~\ref{def:log-SC-algebra} carries a
uniform geometric reading through Pantev--To\"en--Vaqui\'e--Vezzosi
shifted symplectic geometry. The logarithmic $\SCchtop$-algebra structure
encodes a Lagrangian self-intersection
$\mathfrak{S}_b := \mathcal{L} \times_{\mathcal{M}} \mathcal{L}$
as a $(-1)$-shifted symplectic derived stack:
\begin{enumerate}[label=\textup{(\roman*)}]
\item The BV antibracket is the $(-2)$-shifted symplectic
 structure on a derived moduli $\mathcal{M}_{\mathrm{vac}}$
 of vacua; the quantum master equation is the closure condition $d\omega_{-2}=0$.
\item The product decomposition $\omega_k = \omega_k^{\mathrm{hol}} \otimes \omega_k^{\mathrm{top}}$
 is the decomposition of the self-intersection $\mathfrak{S}_b$: the
 $(-1)$-shifted symplectic structure splits into a holomorphic Poisson bracket
 (from $\FM_k(\C)$) and an associative product (from $\Conf_k(\R)$),
 producing the two colors of $\SCchtop$.
\item The Fulton--MacPherson compactification (Theorem~\ref{thm:FM-calculus})
 resolves the collision strata; these strata are the Lagrangian intersection strata.
\end{enumerate}
The Steinberg presentation (Theorem~\ref{thm:steinberg-presentation})
states this directly: the self-intersection data of $\mathcal{L}$ determines the
algebra structure.
\end{remark}

\subsection{From prefactorization to $W(\mathsf{SC}^{\mathrm{ch,top}})$-algebra}
\label{subsec:from-prefact}
\begin{theorem}[Swiss-cheese algebra of observables]
\label{thm:Obs-is-SC}
Under Theorem~\ref{thm:physics-bridge}, the BV observable prefactorization algebra $\mathsf{Obs}$ is a logarithmic $\SCchtop$-algebra (Definition~\ref{def:log-SC-algebra}), functorial under quasi-isomorphisms of HT theories.
\end{theorem}

\begin{proof}
Proposition~\ref{prop:prefact-Obs} verifies the product decomposition; Theorem~\ref{thm:recognition-SC} supplies the $C_\ast(W(\SCchtop))$-algebra structure.
\end{proof}

\subsection{Higher operations via configuration-space integrals}
\label{subsec:mk-integrals}

The $W(\SCchtop)$-algebra structure on $\mathsf{Obs}$ is abstract.
To write the higher products $m_k$ as explicit multilinear operations
on a boundary algebra, fix a small rectangle $D \times I$ and
extract the structure maps from the renormalized time-ordered
product by integrating against configuration-space kernels. The
Stasheff relations for these $m_k$ are Stokes' theorem on
$\FM_k(\C) \times E_1(m)$.

Write $A:=\mathsf{Obs}(D\times I)$ for a small rectangle. The higher multiplications
\[
m_k : A^{\otimes k}\longrightarrow A((\lambda_1))\cdots((\lambda_{k-1}))
\]
of degree $2-k$ are the operations
\(m_k^{\mathrm{BV}}\) of \eqref{eq:3dht-boundary-operations}: integrate
the graph weight \(W_\Gamma(P_{0<L},I)\) over
\(\FM_k(\C)\times E_1(m)\), then take the residue on the boundary face
indexed by \(\Gamma\). The boundary faces yield Stokes' identities
matching the \(A_\infty\) relations.  The bar-side comparison is the
transferred equality \eqref{eq:3dht-bv-bar-transfer-comparison}; the
raw graph residues are not identified with the bar coderivation until
the boundary SDR and comparison maps have been chosen.

\subsection{Well-definedness and sesquilinearity on cohomology}
\label{subsec:well-def-sesqui}

The operations $\eqref{eq:prod-def}$--$\eqref{eq:bracket-def}$ of
Section~\ref{sec:axioms-Ainfty-chiral} descend to cohomology when every algebraic check for a PVA passes: well-definedness
on cohomology classes, sesquilinearity, graded commutativity and
associativity of the product, skewsymmetry and Jacobi for the
bracket, and the Leibniz rule joining them. The four
subsections verify each; the canonical proofs sit in the
PVA descent chapter.

\begin{lemma}[Well-definedness]
\label{lem:well-definedness-prod-bracket}
Definitions \eqref{eq:prod-def} and \eqref{eq:bracket-def} are independent of the choice of representatives of $[a]$ and $[b]$.
\end{lemma}

\begin{proof}
Replace $a\mapsto a+Qx$ (and similarly for $b$). The $n=2$ $A_\infty$ identity (Definition~\ref{def:ainfty_chiral}, Equation~\eqref{eq:ainfty-relation-raw} at $n=2$) reads
\[
 \sum_{\substack{i+j=3 \\ i,j\ge 1}}\;\sum_{s=0}^{2-j}
 (-1)^{\epsilon(s,j)}\;
 m_i\bigl(a_1,\dots,a_s,\,m_j(a_{s+1},\dots,a_{s+j}),\,a_{s+j+1},\dots,a_2\bigr) = 0,
\]
where $\epsilon(s,j) = |a_1|+\cdots+|a_s|$ (Koszul sign convention from Section~\ref{subsec:symmetry}).
Expanding: $(i,j)=(1,2)$ gives $s=0$ only, so $\epsilon(0,2)=0$ and the term is $m_1(m_2(a_1,a_2))$; $(i,j)=(2,1)$ gives $s=0$ with $\epsilon(0,1)=0$ yielding $m_2(m_1(a_1),a_2)$, and $s=1$ with $\epsilon(1,1)=|a_1|$ yielding $(-1)^{|a_1|}m_2(a_1,m_1(a_2))$. Writing $Q=m_1$, the identity becomes
\[
 Q\bigl(m_2(a,b)\bigr) + m_2(Qa,b) + (-1)^{|a|}m_2(a,Qb) = 0.
\]
Thus the first two terms have sign $+1$ and the third has sign $(-1)^{|a|}$, as claimed. Setting $a\mapsto a+Qx$ gives $m_2(a+Qx,b) - m_2(a,b) = m_2(Qx,b) = -Q(m_2(x,b)) - (-1)^{|x|}m_2(x,Qb)$, which is $Q$-exact when $b$ is $Q$-closed. The analogous argument applies when replacing $b$. Together with sesquilinearity \eqref{eq:sesqui-left}--\eqref{eq:sesqui-right}, this shows that the variation of $m_2$ is $Q$-exact, hence the classes of $m_2^{\mathrm{reg}}$ and $m_2^{\mathrm{sing}}$ are invariant. Since $Q = m_1$ acts on the algebra elements and does not involve the spectral parameter $\lambda$, it commutes with the Laurent decomposition into regular and singular parts.
\end{proof}

\begin{proposition}[Sesquilinearity on $H$]
\label{prop:sesqui-H}
The bracket satisfies the vertex sesquilinearity on cohomology:
\[
\{\partial[a]{}_\lambda [b]\} = -\lambda \{[a]{}_\lambda [b]\},\qquad
\{[a]{}_\lambda \partial[b]\} = (\lambda+\partial)\{[a]{}_\lambda [b]\}.
\]
\end{proposition}

\begin{proof}
Apply \eqref{eq:sesqui-left}–\eqref{eq:sesqui-right} to $m_2$, take regular/singular projections, and pass to cohomology.
\end{proof}

\subsection{Commutativity and associativity of the product}
\label{subsec:comm-assoc}
\begin{lemma}[Graded commutativity]
\label{lem:commutative}
$[a]\cdot[b] = (-1)^{|a||b|}[b]\cdot[a]$.
\end{lemma}

\begin{proof}
See Proposition~\ref{prop:product-commutative}; the canonical proof uses the
exchange cylinder construction (Construction~\ref{const:exchange-cylinder}).
\end{proof}

\begin{theorem}[Associativity of the product on cohomology]
\label{thm:assoc}
The product $\cdot$ is associative on $H^\bullet(A,Q)$.
\end{theorem}

\begin{proof}
See Proposition~\ref{prop:product-associative}; the canonical proof uses
the regular--regular projection of the degree-$3$ $A_\infty$ relation.
\end{proof}

\subsection{Skewsymmetry and Jacobi for the bracket}
\label{subsec:skew-Jacobi}
\begin{proposition}[Skewsymmetry]
\label{prop:skewsym}
$\{[a]{}_\lambda [b]\} =
-(-1)^{|a||b|}\{[b]{}_{-\lambda-\partial} [a]\}$.
\end{proposition}

\begin{proof}
See Proposition~\ref{prop:PVA2_proof}; the canonical proof uses the
exchange cylinder in the singular sector plus the Borel transform
$\zeta\mapsto\lambda$.
\end{proof}

\begin{theorem}[Jacobi identity]
\label{thm:Jacobi-PVA}
\label{thm:Jacobi-bv}
For all $[a],[b],[c]\in H^\bullet(A,Q)$,
\[
\{[a]{}_\lambda \{[b]{}_\mu [c]\}\} - (-1)^{(|a|+1)(|b|+1)} \{[b]{}_\mu \{[a]{}_\lambda [c]\}\} = \{\{[a]{}_\lambda [b]\}{}_{\lambda+\mu} [c]\}.
\]
\end{theorem}

\begin{proof}
See Theorem~\ref{thm:Jacobi} in Section~\ref{sec:PVA_descent}; the
canonical proof uses the full three-face singular--singular Stokes
argument on $\FM_3(\C)$ with Arnold--Orlik--Solomon corner cancellation.
\end{proof}

\subsection{Leibniz rule}
\label{subsec:Leibniz}
\begin{theorem}[Leibniz rule]
\label{thm:Leibniz-PVA}
\label{thm:Leibniz-bv}
For all $[a],[b],[c]\in H^\bullet(A,Q)$,
\[
\{[a]{}_\lambda ([b]\cdot [c])\} = \{[a]{}_\lambda [b]\}\cdot [c] + (-1)^{(|a|+1)|b|}\,[b]\cdot \{[a]{}_\lambda [c]\}.
\]
\end{theorem}

\begin{proof}
See Proposition~\ref{prop:PVA3_proof} in Section~\ref{sec:PVA_descent};
the canonical proof uses the mixed regular--singular three-face Stokes
projection on $\FM_3(\C)$.
\end{proof}

\begin{corollary}[PVA on cohomology]
\label{cor:PVA-cohomology-bv}
\label{cor:PVA-structure-bv}
The operations \eqref{eq:prod-def}--\eqref{eq:bracket-def} define a
$(-1)$-shifted Poisson vertex algebra structure on $H^\bullet(A,Q)$.
The visible skewsymmetry of the $\lambda$-bracket is the standard
vertex-algebra identity above. The shifted Koszul signs occur in the
Jacobi and Leibniz identities and come from the bar desuspension:
the operations $m_k$ are components of a degree-$+1$ coderivation on
$T^c(s^{-1}\bar A)$.
\end{corollary}

\begin{remark}[Transition to FM calculus]
\label{rem:transition-to-FM}
The $m_k$ are defined; the Stasheff identities are Stokes' theorem on
$\FM_k(\C)$. Codimension-$1$ faces encode compositions
$m_i \circ m_j$; codimension-$2$ corners cancel through the
Arnold--Orlik--Solomon relations.
Section~\ref{sec:FM_calculus} carries this out.
\end{remark}
